% !TeX root = ../thuthesis-example.tex

\chapter{总结与展望}

\section{总结}

本文完成了“基于Spring Boot的分布式推荐系统”的课程大作业报告撰写。系统以类Telegram社交应用为业务背景，将后端设计为Spring Boot分层架构，前端负责页面展示和行为采集，Supabase PostgreSQL负责云端数据持久化，Redis负责缓存与时间线加速，Apifox负责接口联调。

在系统设计方面，本文按照课程模板完成了需求分析、可行性分析、系统架构设计、模块设计、数据库设计和接口设计。系统整体分为客户端层、后端服务层、数据与缓存层和推荐算法层。后端模块包括认证与权限、用户与群组、新闻内容、推荐Feed、行为埋点和实验配置。该设计避免了将所有逻辑堆叠在单一接口中，体现了分布式应用的模块边界意识。

在推荐算法方面，本文重点描述了多阶段推荐流水线。系统先构建请求上下文，再补全用户画像和实验信息，随后通过关注关系、社交图谱、新闻向量、作者兴趣、热门内容和冷启动内容进行多路召回。候选内容进入过滤阶段后，会去除重复、自身内容、已曝光内容、屏蔽用户和敏感词。评分阶段综合行为权重、作者亲和度、内容质量、新鲜度和负反馈惩罚，最终通过TopK选择生成个性化Feed。

在测试方面，本文使用Apifox完成推荐Feed接口的请求配置截图，并通过Supabase MCP读取真实云端数据库结构。数据库中已存在用户、群组、联系人、新闻内容、用户行为、用户兴趣向量和实验配置等表，能够支撑推荐系统的基本数据需求。

\section{展望}

虽然本项目已经形成完整的课程报告和系统设计，但距离真实生产级推荐系统仍有改进空间。

首先，后端可以从单体Spring Boot多模块架构逐步演进为微服务架构。推荐服务、用户服务、内容服务、实验服务和日志服务可以独立部署，通过服务注册、网关和熔断机制提高系统可维护性。

其次，推荐算法可以从规则加权模型升级为可训练模型。当前设计重点强调可解释性和工程可落地性，后续可以引入双塔模型、矩阵分解、图神经网络或Learning to Rank模型，根据真实行为数据自动学习权重。

第三，行为数据闭环仍需加强。后续可以建立离线特征计算任务，定期根据点击、停留、分享和负反馈更新用户兴趣向量，并将模型指标接入实验平台。

第四，安全策略需要完善。Supabase当前存在公开表未启用RLS的问题，真实上线前应启用行级安全策略，补充按用户身份、管理员角色和服务端密钥区分的访问规则，同时对密码字段进行哈希存储和脱敏处理。

最后，测试体系可以进一步自动化。Apifox当前主要用于手动联调和截图记录，后续可增加响应断言、环境变量、前置脚本和自动化测试集合，使接口测试能够随着后端迭代持续执行。
